\chapter{Machten}\label{ch:g8:powers}

Vouw een blad papier $10$ keer dubbel (als je kunt!): de dikte verdubbelt elke keer, en
$10$ verdubbelingen vermenigvuldigen haar met $2^{10} = 1024$. Machten zijn de korte
schrijfwijze van het herhaald vermenigvuldigen; dit hoofdstuk zet de notatie en haar
regels op, met een bijzondere rol voor de machten van $10$.

\section{Definitie}

\begin{definition}[Macht]\label{def:g8:powers:def}
Voor een getal $a$ en een heel getal $n \geq 1$:
\[
a^n = \underbrace{a \times a \times \dots \times a}_{n
\text{ factoren}},
\]
gelezen als ``$a$ tot de macht $n$''; $a$ is het \emph{grondtal}, $n$ de
\emph{exponent}\index{exponent}. Bijzondere namen: $a^2$ is ``$a$ kwadraat'', $a^3$
``$a$ tot de derde''. Per afspraak:
\[
a^1 = a,
\qquad
a^0 = 1 \ (a \neq 0),
\qquad
a^{-n} = \frac{1}{a^n}.
\]
\end{definition}

\begin{example}\label{ex:g8:powers:def}
$3^4 = 3 \times 3 \times 3 \times 3 = 81$; \quad $10^3 = 1000$; \quad
$5^{-2} = \frac{1}{25} = 0.04$; \quad $(-2)^3 = -8$ en $(-2)^4 = +16$
(\cref{ex:g8:negprod:powers}). Pas op: $a^n$ is \emph{niet} $a \times n$: \
$2^5 = 32$, en niet $10$.
\end{example}

\begin{omfigure}
\begin{tikzpicture}[scale=0.9]
  \foreach \n/\x in {0/0, 1/1.4, 2/2.8, 3/4.2, 4/5.6} {
    \pgfmathsetmacro{\h}{0.28*pow(2,\n)}
    \draw[omDef, fill=omDef!20] (\x,0) rectangle (\x+0.9,\h);
    \node[below, font=\small] at (\x+0.45,0) {$2^\n$};
    \pgfmathtruncatemacro{\v}{pow(2,\n)}
    \node[above, font=\small] at (\x+0.45,\h) {$\v$};
  }
\end{tikzpicture}

{\small Machten van $2$: elke staaf is twee keer de vorige. Groei door herhaald
vermenigvuldigen loopt veel sneller weg dan groei door herhaald optellen.}
\end{omfigure}

\section{De rekenregels voor machten}

\begin{theorem}[Rekenregels voor machten]\label{thm:g8:powers:rules}
Voor een grondtal $a$ dat niet \omterm{def:g1:counting:zero}{nul} is en hele \omterm{def:g8:powers:def}{exponenten} $m$ en $n$:
\[
a^m \times a^n = a^{m+n},
\qquad
\frac{a^m}{a^n} = a^{m-n},
\qquad
\left(a^m\right)^n = a^{m \times n},
\qquad
(ab)^n = a^n b^n .
\]
\end{theorem}

\begin{proof}[Bewijs door factoren te tellen]
$a^m \times a^n$ zet $m$ factoren $a$ op een rij, gevolgd door nog $n$: samen
$m + n$. $\left(a^m\right)^n$ herhaalt een blok van $m$ factoren $n$ keer: $mn$
factoren. $(ab)^n$ bevat $n$ letters $a$ en $n$ letters $b$, die je kunt hergroeperen.
De regel voor het \omterm{def:g3:division:remainder}{quotiënt} komt van het wegstrepen van $n$ van de $m$ factoren; met de
afspraken $a^0 = 1$ en $a^{-n} = \frac{1}{a^n}$ blijft ze zelfs geldig als
$n \geq m$.
\end{proof}

\begin{example}\label{ex:g8:powers:rules}
\[
2^3 \times 2^5 = 2^8 = 256,
\qquad
\frac{7^6}{7^4} = 7^2 = 49,
\qquad
\left(5^2\right)^3 = 5^6,
\qquad
2^4 \times 5^4 = 10^4 .
\]
Een valstrik: de regels gelden bij een \emph{gemeenschappelijk grondtal} (of, bij de
laatste, een gemeenschappelijke \omterm{def:g8:powers:def}{exponent}). Geen enkele regel vereenvoudigt
$2^3 \times 5^2$ --- reken het gewoon uit: $8 \times 25 = 200$.
\end{example}

\section{Machten van tien}

\begin{proposition}[Machten van tien]\label{prop:g8:powers:ten}
Voor $n \geq 1$ is $10^n = 1\underbrace{0\dots0}_{n}$ en
$10^{-n} = 0.\underbrace{0\dots0}_{n-1}1$. De rekenregels voor machten zeggen: bij het
vermenigvuldigen van machten van tien tel je de \omterm{def:g8:powers:def}{exponenten} op.
\end{proposition}

\begin{example}\label{ex:g8:powers:ten}
$10^4 \times 10^3 = 10^7$; \quad $\dfrac{10^2}{10^5} = 10^{-3} = 0.001$. Grote en
kleine hoeveelheden worden leesbaar:
\[
\text{een miljard} = 10^9,
\qquad
\text{een miljoenste} = 10^{-6}.
\]
Samen met decimalen: $3.2 \times 10^5 = 320\,000$ en
$4.7 \times 10^{-3} = 0.0047$ (schuif het \omterm{def:g4:decimals:point}{decimaalteken},
\cref{prop:g6:decimals:shift}). Het stelselmatige gebruik van deze schrijfwijze ---
de \emph{wetenschappelijke notatie} --- wordt in \cref{ch:g9:fractions} uitgewerkt.
\end{example}

\begin{example}[Grootteordes]\label{ex:g8:powers:magnitude}
Licht legt ongeveer $3 \times 10^8$ m/s af; een jaar heeft ongeveer
$3.2 \times 10^7$ seconden. Een lichtjaar is dus ongeveer
\[
3 \times 3.2 \times 10^{8+7} \approx 10 \times 10^{15} = 10^{16}
\text{ m}
\]
--- tien miljoen miljard meter. Met machten van tien passen astronomische berekeningen
op één regel.
\end{example}

\begin{method}[Een uitdrukking met machten vereenvoudigen]\label{met:g8:powers:simplify}
\begin{enumerate}
  \item Groepeer de factoren per grondtal;
  \item pas binnen elk grondtal de rekenregels voor de \omterm{def:g8:powers:def}{exponenten} toe;
  \item reken de kleine machten die overblijven uit, of laat het antwoord als macht
        staan als het groot is.
\end{enumerate}
\end{method}

\begin{example}\label{ex:g8:powers:simplify}
\[
\frac{3^5 \times 3^{-2} \times 4^2}{3^2}
= \frac{3^{5 + (-2)}}{3^2} \times 16
= 3^{3 - 2} \times 16
= 3 \times 16 = 48 .
\]
\end{example}

\section{Oefeningen}

\begin{exercise}[$\star$]\label{exo:g8:powers:1}
Reken uit:
\[
2^6, \qquad 3^3, \qquad 10^5, \qquad 1^{100}, \qquad 0.1^2, \qquad
6^0 .
\]
\end{exercise}

\begin{exercise}[$\star$]\label{exo:g8:powers:2}
Reken uit:
\[
(-3)^2, \qquad -3^2, \qquad (-1)^{15}, \qquad (-5)^3, \qquad
\left(\tfrac{2}{3}\right)^2 .
\]
\end{exercise}

\begin{exercise}[$\star$]\label{exo:g8:powers:3}
Schrijf als één macht:
\[
7^4 \times 7^5, \qquad
\frac{2^9}{2^3}, \qquad
\left(10^3\right)^4, \qquad
5^6 \times 5^{-2}, \qquad
3^4 \times 7^4 .
\]
\end{exercise}

\begin{exercise}[$\star$]\label{exo:g8:powers:4}
Schrijf als \omterm{def:g6:decimals:places}{decimaal getal}: $10^{-2}$; \quad $4 \times 10^3$; \quad
$2.5 \times 10^{-4}$; \quad $10^0$.
\end{exercise}

\begin{exercise}[$\star$]\label{exo:g8:powers:5}
Schrijf met een macht van tien: honderdduizend; een tiende; tien miljard;
$0.000\,001$.
\end{exercise}

\begin{exercise}[$\star$]\label{exo:g8:powers:6}
Vereenvoudig en reken daarna uit:
\[
\frac{10^7 \times 10^{-3}}{10^2},
\qquad
\frac{2^5 \times 2^4}{2^6},
\qquad
\left(2^2\right)^3 \times 2^{-4} .
\]
\end{exercise}

\begin{exercise}[$\star$]\label{exo:g8:powers:7}
Waar of niet waar? Verbeter de onjuiste.
\[
2^3 \times 2^4 = 2^{12}; \qquad
5^2 + 5^3 = 5^5; \qquad
(3^2)^4 = 3^8; \qquad
10^3 \times 10^3 = 100^3 .
\]
\end{exercise}

\begin{exercise}[$\star\star$]\label{exo:g8:powers:8}
Een gerucht doet de ronde: op dag 1 kennen drie mensen het; elke dag vertelt elke
persoon die het kent, het aan drie nieuwe mensen. Schrijf met een macht het aantal
\emph{nieuwe} mensen dat op dag $4$ wordt ingelicht, en reken uit hoeveel mensen het
gerucht aan het einde van dag 4 kennen (de oorspronkelijke drie meegerekend).
\end{exercise}

\begin{exercise}[$\star\star$]\label{exo:g8:powers:9}
Een blad papier is $0.1$ mm dik, dus $10^{-4}$ m. Elke keer dubbelvouwen verdubbelt de
dikte.
\begin{enumerate}
  \item Druk de dikte na $10$ vouwen als een \omterm{def:g2:mult:def}{product} uit, en reken haar in centimeters
        uit ($2^{10} = 1024$).
  \item Na $42$ vouwen zou de dikte $2^{42} \times 10^{-4}$ m zijn, met
        $2^{42} \approx 4.4 \times 10^{12}$. Laat zien dat dat meer is dan de afstand
        van de aarde tot de maan, ongeveer $3.8 \times 10^8$ m.
\end{enumerate}
\end{exercise}

\begin{exercise}[$\star\star$]\label{exo:g8:powers:10}
Zet op volgorde, van klein naar groot, zonder rekenmachine:
\[
2^{10}, \qquad 10^3, \qquad 3^6, \qquad 5^4 .
\]
(Reken ze elk uit; $2^{10}$ en $10^3$ zijn beroemde buren.)
\end{exercise}

\begin{exercise}[$\star\star\star$]\label{exo:g8:powers:11}
Wat is groter, $2^{100}$ of $10^{30}$? Gebruik $2^{10} = 1024 > 10^3$ om
$2^{100} = \left(2^{10}\right)^{10}$ met $\left(10^3\right)^{10}$ te vergelijken.
\end{exercise}

\section{Opgave: het schaakbord en de machten van twee}

\begin{problem}[{Weekendopgave --- de meetkundige \omterm{def:g1:addition:def}{som}
$1 + 2 + 4 + \dots + 2^{n-1} = 2^n - 1$, van een beroemde legende tot binaire
getallen}]
\label{pb:g8:powers:1}
De legende: als belofte voor het uitvinden van het schaakspel vroeg de wijze Sissa zijn
koning één graankorrel op het eerste veld van het bord, twee op het tweede, vier op het
derde --- verdubbelend van veld tot veld, tot het vierenzestigste. De koning lachte om
zoveel bescheidenheid. In deze opgave reken je uit wat de koning beloofde, met de
rekenregels voor machten van \cref{thm:g8:powers:rules}, en eindig je waar het verhaal
in stilte naartoe leidt: de binaire getallen in elke computer.

\medskip
\noindent\textbf{Deel I --- De verdubbeltruc.}
Zij voor $n \geq 1$ $S_n$ het totale aantal korrels op de eerste $n$ velden.
\begin{enumerate}
  \item Druk het aantal korrels op veld $k$ uit als een macht van $2$. Welke macht
        staat op veld $64$?
  \item Reken $S_1$, $S_2$, $S_3$, $S_4$ en $S_5$ uit, en vergelijk elk met een
        nabijgelegen macht van $2$. Vermoed een formule voor $S_n$.
  \item De \emph{verdubbeltruc}: schrijf de sommen $S_n$ en $2 \times S_n$ onder
        elkaar, trek ze af, en bewijs je vermoeden:
        \[
        S_n = 1 + 2 + 4 + \dots + 2^{n-1} = 2^n - 1 .
        \]
  \item Hoeveel korrels beloofde de koning in totaal? Druk het antwoord uit met een
        macht van $2$, en maak de klassieke opmerking af: ``het hele bord bevat één
        korrel minder dan één enkel vijfenzestigste veld zou bevatten.''
  \item Laat zien dat de tweede \omterm{def:g3:division:half}{helft} van het bord (de velden $33$ tot $64$)
        \emph{precies} $2^{32}$ keer zoveel korrels bevat als de eerste \omterm{def:g3:division:half}{helft}.
\end{enumerate}

\medskip
\noindent\textbf{Deel II --- Hoe groot is $2^{64}$?}
De vergelijking $2^{10} = 1024 > 10^3$ uit \cref{exo:g8:powers:11} is de sleutel tot
alle schattingen hieronder.
\begin{enumerate}[resume]
  \item Laat zien dat $2^{64} = 2^4 \times \left(2^{10}\right)^6 > 1.6 \times
        10^{19}$.
  \item Een graankorrel weegt ongeveer $0.05$ g, dus $5 \times 10^{-2}$ g. Laat zien
        dat het beloofde graan meer dan $8 \times 10^{17}$ g weegt, en zet dat om in
        tonnen ($1$ ton $= 10^6$ g).
  \item De hele wereld oogst momenteel ongeveer $8 \times 10^8$ ton graan per jaar.
        Hoeveel jaar wereldoogst beloofde de koning minstens?
  \item Zoek het kleinste hele getal $n$ waarvoor $2^n > 10^6$ --- dus hoeveel
        verdubbelingen er nodig zijn om een miljoen voorbij te gaan. (Reken $2^{19}$ en
        $2^{20}$ precies uit met $2^{10} = 1024$.)
  \item Een grappenmaker biedt je een maandsalaris aan: $1$ cent op dag $1$, en daarna
        elke dag het dubbele van de dag ervoor. Op welke dag gaat het \emph{dagloon}
        alleen al voor het eerst boven een miljoen euro ($10^8$ cent)? (Reken $2^{26}$
        en $2^{27}$ precies uit.)
\end{enumerate}

\medskip
\noindent\textbf{Deel III --- Binaire gewichten.}
Een marktkoopvrouw heeft vijf gewichten: $1$, $2$, $4$, $8$ en $16$ gram, van elk één.
Ze legt er een aantal van op de ene schaal van een balans om waren op de andere te
wegen.
\begin{enumerate}[resume]
  \item Welke gewichten legt ze neer om $21$ g te wegen? En om $27$ g te wegen?
  \item Leg uit waarom elk doelgewicht van $16$ g of meer het gewicht van $16$ g
        \emph{moet} gebruiken, en waarom elk doelgewicht van $15$ g of minder het
        \emph{niet} mag gebruiken. (Vraag 3 vertelt je wat de gewichten $1, 2, 4, 8$
        hoogstens kunnen bereiken.) Leg uit waarom dezelfde redenering zich bij het
        volgende grootste gewicht herhaalt, en zo bij elke stap.
  \item Leid af dat elk heel doelgewicht van $1$ tot $31$ g gewogen kan worden, en
        \emph{op precies één manier}: elk getal tussen $1$ en $31$ is op één enkele
        manier een \omterm{def:g1:addition:def}{som} van verschillende machten van $2$.
  \item De koopvrouw koopt een zesde gewicht, van $32$ g. Tot welk doelgewicht kan ze
        nu wegen? Schrijf $45$ g als een \omterm{def:g1:addition:def}{som} van verschillende machten van $2$.
  \item In een computer wordt een ``64-bits'' getal op $64$ velden bewaard, die elk een
        $0$ of een $1$ bevatten --- veld $k$ draagt $2^{k-1}$ bij als het een $1$
        bevat, zoals de korrels van de legende. Leg met deel I uit waarom de hele
        getallen die zo'n machine kan bewaren precies van $0$ tot $2^{64} - 1$ lopen.
\end{enumerate}
\end{problem}
